% Standalone appendix section. Requires amsmath, amssymb, amsthm and
% a proposition environment. No custom manuscript macros are used.
\section{Conditional estimation and recovery guarantees}
\label{app:theory}

The results below describe the conditions under which a calibrated residual can
support execution recovery. The metric concerns robot dynamics and their
low-level controller under a shared supplied command stream. It is not a
metric for the VLA feedback loop, and the reported robot experiments do not
verify this physical metric. In particular,
the estimator results and the trajectory results have different assumptions.
All norms are Euclidean or their induced operator norms unless stated otherwise.

\subsection{Residual history, calibration error, and identifiability}

Let $u_k$ be the actual command sent by the adapter before the additive fault.
An idealized local FIR description is
\begin{equation}
 y_k=\sum_{\ell=0}^{K}G_\ell(u_{k-\ell}+f_{k-\ell})+c+b_y+\varepsilon_k^{\rm FIR},
 \qquad
 \widehat y_k=\sum_{\ell=0}^{K}H_\ell u_{k-\ell}+\widehat c.
 \label{eq:app-fir}
\end{equation}
Here $b_y$ is a constant healthy residual bias in measurement units, and
$\varepsilon_k^{\rm FIR}$ includes departures from the local FIR description. Define
$G_\Sigma=\sum_\ell G_\ell$, a calibrated nonsingular sensitivity $S$, a
healthy bias calibration $\widehat b_y=S b_{\mathrm{cal}}$, and
\begin{align}
 z_k&=S^{-1}(y_k-\widehat y_k-\widehat b_y)=f_k+w_k,
 \label{eq:app-observation}\\
 w_k&=S^{-1}\left[
  \sum_{\ell=0}^{K}G_\ell(f_{k-\ell}-f_k)
  +(G_\Sigma-S)f_k+(b_y-\widehat b_y)+n_k\right],
 \label{eq:app-observation-error}\\
 n_k&=\varepsilon_k^{\rm FIR}+\sum_{\ell=0}^{K}(G_\ell-H_\ell)u_{k-\ell}
       +(c-\widehat c).\nonumber
\end{align}
Thus $e_k^{\rm res}-\widehat b_y\simeq Sf_k$ is a steady or slowly varying approximation, not an
instantaneous identity at fault onset. If $\|n_k\|\le\bar n$,
$\|b_y-\widehat b_y\|\le\Delta_b$,
$\|G_\Sigma-S\|\le\Delta_S$, $\|f_k\|\le F$, and
$\|f_{k+1}-f_k\|\le\nu$, then
\begin{equation}
 \|w_k\|\le\epsilon:={\|S^{-1}\|}
 \left(\bar n+\Delta_b+\Delta_S F+
       \nu\sum_{\ell=0}^{K}\ell\|G_\ell\|\right).
 \label{eq:app-fir-bound}
\end{equation}
This follows by $\|f_{k-\ell}-f_k\|\le\ell\nu$ and the triangle
inequality. At an abrupt change, the exact history term in
\eqref{eq:app-observation-error} is preferable to this uniform bound; for a
constant post-onset fault it vanishes after the FIR window has filled. An
incorrectly initialized history contributes to $n_k$ until that history has
been replaced. Bounds on predictor error and sensitivity mismatch must hold
on the corrected trajectories, not only on healthy calibration data.

\paragraph{Proof of Proposition~\ref{prop:ambiguity} and its scope.}
In the constant-sensitivity model $e_k^{\rm res}=Sf+b_y+\varepsilon_k^{\rm obs}$, with constant unknown
$f$ and unrestricted constant unknown $b_y$, $f$ and $b_y$ cannot be separately
identified. For every vector $v$, the pair $(f+v,b_y-Sv)$ gives exactly the
same observations. With known $b_y$ and nonsingular $S$, the noiseless residual
does determine $f$. The identity $S(f+v)+(b_y-Sv)=Sf+b_y$ proves the ambiguity. For known bias,
$f=S^{-1}(e_k^{\rm res}-b_y)$ in the noiseless model.

A known healthy pre-onset interval can therefore provide the bias calibration
within the same episode, subject to noise and history transients. More generally,
a known fault subspace or time-varying sensitivity can change identifiability;
the proposition is not an impossibility statement for every sensor and model.
An offline mean of the raw healthy residual estimates $b_y$ only when the
remaining healthy modeling error is zero-mean under that sampling distribution.
Otherwise it estimates an aggregate healthy offset whose transfer to corrected
trajectories needs its own error bound. A mean of the settled, attenuated legacy
estimate is a different quantity and must not be silently substituted for this
raw residual calibration. For example, if a constant healthy observation in
fault units is $b_0$ and attenuation is fixed at $s_0$, a settled legacy
estimate equals $s_0b_0$ when projection is inactive. Subtracting that estimate
before attenuation leaves $s_0(b_0-s_0b_0)=s_0(1-s_0)b_0$, generally nonzero.
The averaging quantity, units, and subtraction order therefore affect the
identification error and must match the executed implementation.

\subsection{The two additive update laws and the physical command clock}

Let $C\subset\mathbb R^d$ be closed and convex, assume $f_k\in C$, and let
$\Pi_C$ denote Euclidean projection. Let $\widehat f_0\in C$ and
$E_k=\|f_k-\widehat f_k\|$. The box constraint in the experiments is one
choice of $C$; it bounds the estimate but does not by itself guarantee safe
plant motion.

\paragraph{Full-error counterpart of Proposition~\ref{thm:estimator}.}
Suppose \eqref{eq:app-observation} holds with $\|w_k\|\le\epsilon_k$ and
$\|f_{k+1}-f_k\|\le \nu_k$. First consider one active innovation update
\begin{equation}
 \widehat f_{k+1}=\Pi_C\{\widehat f_k+\alpha_k(z_k-\widehat f_k)\},
 \qquad 0\le\alpha_k\le1.
 \label{eq:app-innovation}
\end{equation}
Then
\begin{equation}
 E_{k+1}\le(1-\alpha_k)E_k+\alpha_k\epsilon_k+\nu_k.
 \label{eq:app-estimator-recursion}
\end{equation}
Writing $P_{j,k}=\prod_{i=j}^{k-1}(1-\alpha_i)$ with empty products equal to
one gives the finite-horizon bound
\begin{equation}
 E_k\le P_{0,k}E_0+
 \sum_{j=0}^{k-1}P_{j+1,k}(\alpha_j\epsilon_j+\nu_j).
 \label{eq:app-estimator-product}
\end{equation}
If $\epsilon_k\le\epsilon$, $\nu_k\le\nu$, and
$\alpha_k\ge\underline\alpha>0$, then, with $q=1-\underline\alpha$,
\begin{equation}
 E_k\le q^k E_0+\epsilon+\frac{\nu}{\underline\alpha}.
 \label{eq:app-estimator-uniform}
\end{equation}
In particular, an exactly modeled constant fault converges exponentially.
\begin{proof}
Since $f_k\in C$, projection nonexpansiveness gives
\begin{align*}
 \|\widehat f_{k+1}-f_k\|
 &\le\|(1-\alpha_k)(\widehat f_k-f_k)+\alpha_k w_k\|\\
 &\le(1-\alpha_k)E_k+\alpha_k\epsilon_k.
\end{align*}
Adding $\|f_{k+1}-f_k\|$ proves the recursion; induction proves the product
bound. Also $P_{0,k}\le q^k$,
$\sum_{j=0}^{k-1}\alpha_jP_{j+1,k}=1-P_{0,k}\le1$, and
$\sum_{j=0}^{k-1}P_{j+1,k}\le\sum_{i=0}^{k-1}q^i
\le1/\underline\alpha$, proving the uniform bound.
\end{proof}

The innovation normalizer corresponds to
\begin{equation}
 \alpha_k=\frac{\gamma}{1+\|P_rS(z_k-\widehat f_k)\|^2/\rho^2},
 \qquad 0<\gamma\le1.
 \label{eq:app-alpha}
\end{equation}
If $C$ has diameter $D_C$ and $\|w_k\|\le\epsilon$, then
$\underline\alpha=\gamma/[1+\|P_rS\|^2(D_C+\epsilon)^2/\rho^2]$ is a
valid, possibly conservative lower bound. This estimator does not require
online persistent excitation when the full-rank, constant $S$ is already
known: the fault is directly observed through $S$. Estimating $S$ from healthy
data, and estimating multiplicative faults, are separate identification tasks.

\paragraph{Deadzone placement.}
Equation~\eqref{eq:app-estimator-product} remains valid when a gate sets
$\alpha_k=0$. The uniform exponential bound then requires a positive lower
bound on accumulated effective gain over every fixed-length sliding window,
not merely a positive ungated gain or some active updates.
For a full-channel innovation gate that skips updates when
$\|S(z_k-\widehat f_k)\|\le\delta$, an inactive step implies
\begin{equation}
 E_k\le\epsilon+\|S^{-1}\|\delta.
 \label{eq:app-deadzone}
\end{equation}
This full-error implication requires the full invertible $S$. A selected gate
$\|P_rS(z_k-\widehat f_k)\|\le\delta$ gives it only when $P_rS$ is
injective on the estimated fault space, with $\|S^{-1}\|$ replaced by
$1/\sigma_{\min}(P_rS)$. A rank-deficient selection cannot bound the
unobserved error components. For constant $f$ and bounded noise, this ball is forward invariant under
active steps, and distance above it decreases by at least the ungated factor
$q$ on every active step. Indeed, the inactive-step bound follows directly
from $z_k-\widehat f_k=f_k-\widehat f_k+w_k$; for active steps,
\eqref{eq:app-estimator-recursion} maps that ball into itself and contracts
any excess. A raw-residual gate has no corresponding error implication:
after a fault disappears, $e_k^{\rm res}$ can be small while $\widehat f_k$ is not.
The distinct implementation that zeroes the observation but retains the
$-\widehat f_k$ leakage decays on inactive steps rather than freezing. These
gate conventions must be reported separately.

\begin{proposition}[Attenuation in the legacy update]
\label{prop:app-legacy}
For
\begin{equation}
 \widehat f_{k+1}=\Pi_C\{(1-\gamma)\widehat f_k+\gamma s_k z_k\},
 \qquad 0<\gamma\le1,\quad 0\le s_k\le1,
 \label{eq:app-legacy}
\end{equation}
the estimation error obeys
\begin{equation}
 E_{k+1}\le(1-\gamma)E_k+
 \gamma\{(1-s_k)\|f_k\|+s_k\epsilon_k\}+\nu_k.
 \label{eq:app-legacy-bound}
\end{equation}
For constant noiseless observations and constant $s$, the unique fixed point
is $\Pi_C(sf)$, approached with contraction factor at most $1-\gamma$.
If $sf\in C$, the fixed point is $sf$, which is biased unless $s=1$ or $f=0$.
\end{proposition}
\begin{proof}
Apply projection nonexpansiveness relative to $f_k$ and use
$s_kz_k-f_k=-(1-s_k)f_k+s_kw_k$ to obtain the bound. For constant target
$sf$, the update map is a contraction with factor $1-\gamma$.
Let $p=\Pi_C(sf)$. The characterization
$\langle sf-p,x-p\rangle\le0$ for all $x\in C$ also implies
$\Pi_C\{p+\gamma(sf-p)\}=p$. Thus $p$ is its unique fixed point.
\end{proof}

The documented observation-attenuated normalization has
\[
s_k=\frac{\mathbf{1}[\|P_re_k^{\rm res}\|>\delta]}
{1+\|P_re_k^{\rm res}\|^2/\rho^2},
\]
subject to the selected-channel norm and calibration convention. Its guaranteed error
includes attenuation. The innovation proposition does not retrospectively
give an unbiased-estimation guarantee to experiments run with this law.


\paragraph{Update and hold are distinct from target-only gating.}
Let $k$ index every physical command, including steps with no estimator update.
The associated measurement becomes available after executing $u_k$ and can
first affect $\widehat f_{k+1}$. A full-update indicator $\chi_k\in\{0,1\}$
gives the exact scheduled maps
\begin{align}
\widehat f_{k+1}
 &=\Pi_C\{(1-\chi_k\gamma)\widehat f_k+
                       \chi_k\gamma s_kz_k\},\label{eq:app-scheduled-legacy}\\
\widehat f_{k+1}
 &=\Pi_C\{\widehat f_k+\chi_k\alpha_k(z_k-\widehat f_k)\}.
\label{eq:app-scheduled-innovation}
\end{align}
Because $\widehat f_k\in C$, $\chi_k=0$ holds the estimate exactly.
Applying the two projection arguments above with the effective step sizes
$\chi_k\gamma$ or $\chi_k\alpha_k$ proves
the full-error counterparts of Eqs.~\eqref{eq:legacybound}--\eqref{eq:estbound}. In particular, a full hold
has $E_{k+1}\le E_k+\nu_k$. In contrast, zeroing a legacy observation while
retaining leakage uses $\chi_k=1,s_k=0$ and gives
$E_{k+1}\le(1-\gamma)E_k+\gamma\|f_k\|+\nu_k$.
This can move the estimate toward zero even when the physical fault persists.
The schedule indicator must represent the executed code path; it is not
inferred from the displayed deadzone threshold.

Both recursions have the form
\begin{equation}
 B_{k+1}=(1-\beta_k)B_k+\beta_k\omega_k+\nu_k,\qquad B_0\ge E_0,
 \label{eq:app-schedule-envelope}
\end{equation}
where $(\beta_k,\omega_k)$ equals
$(\chi_k\gamma,(1-s_k)\|f_k\|+s_k\epsilon_k)$ for observation attenuation
and $(\chi_k\alpha_k,\epsilon_k)$ for innovation adaptation.
With $P_{j,k}^{\beta}=\prod_{i=j}^{k-1}(1-\beta_i)$,
\begin{equation}
 E_k\le B_k=P_{0,k}^{\beta}B_0+
 \sum_{j=0}^{k-1}P_{j+1,k}^{\beta}(\beta_j\omega_j+\nu_j).
 \label{eq:app-schedule-product}
\end{equation}
This follows by induction and is valid for arbitrary finite schedules,
including long holds and a fault jump represented by a single large $\nu_j$.
For constant $s$ and noiseless constant $f$, the distance to the legacy
fixed point $\Pi_C(sf)$ is at most $(1-\gamma)^{N_k}$ times its initial
value, where $N_k=\sum_{j=0}^{k-1}\chi_j$. It need not decay
at the same rate per physical step. Converting the count into seconds requires
actual timestamps; a positive ungated gain alone supplies no such rate.

\paragraph{Selected-coordinate proof of Proposition~\ref{thm:estimator}.}
For a box $\mathcal C=\prod_i[\ell_i,u_i]$ and binary diagonal $D$,
coordinatewise clipping gives, for any $v$ and $f\in\mathcal C$,
\begin{equation}
\|D(\Pi_{\mathcal C}(v)-f)\|^2
=\sum_{i:D_{ii}=1}|\Pi_{[\ell_i,u_i]}(v_i)-f_i|^2
\le\|D(v-f)\|^2.
\label{eq:masked-projection}
\end{equation}
Apply this inequality to each scheduled update, expand $z_k=f_k+w_k$,
and add $\|D(f_{k+1}-f_k)\|$. This proves the two main masked recursions.
The scalar envelope and product in Eqs.~\eqref{eq:app-schedule-envelope}--
\eqref{eq:app-schedule-product} then hold with $B_k^D,E_k^D,\nu_k^D$ and
$\omega_k^D=(1-s_k)\|Df_k\|+s_k\epsilon_k^D$ for attenuation, or
$\omega_k^D=\epsilon_k^D$ for innovation.

This conclusion requires separable projection. For a counterexample with a
coupled set, take the unit disk, $f=(1/2,0)$, $v=(1/2,2)$, and
$D=\operatorname{diag}(1,0)$. Then $D(v-f)=0$, whereas the first coordinate
of the projection is $1/\sqrt{17}\ne1/2$. Thus selected error can increase
through projection even when the selected input error is zero. The full-error
proof for arbitrary closed convex sets remains valid.

\paragraph{Application snapshots and buffered commands.}
Let $\widehat f_k^{\rm app}=\widehat f_{\tau_k}$ be the estimate used in
physical command $k$, $0\le\tau_k\le k$. Define
$E_k^{D,\rm app}=\|D(\widehat f_{\tau_k}-f_k)\|$.
Inserting respectively $\widehat f_k$ and $f_{\tau_k}$ gives
\begin{align}
E_k^{D,\rm app}&\le E_k^D+
\sum_{j=\tau_k}^{k-1}\|D(\widehat f_{j+1}-\widehat f_j)\|,
\label{eq:app-delay-increments}\\
E_k^{D,\rm app}&\le E_{\tau_k}^D+\sum_{j=\tau_k}^{k-1}\nu_j^D.
\label{eq:app-delay-age}
\end{align}
Taking the minimum of these valid bounds proves Eq.~\eqref{eq:applied-delay}.
For maximum age $d$ and masked drift $\nu$, the second is
$B_{\tau_k}^D+d\nu$. An episode-long hold uses its full age. No delay
limit is needed for the finite-horizon forms. Setting $D=I$ yields the
full-error bounds $E_k^{\rm app}\le B_k^{\rm app}$ used for a general
convex projection or nonbinary correction map. Snapshot creation time, full
holds, and target-only gates are distinct quantities. If one decoder chunk
uses multiple estimates, its commands must be represented at the matching
finer resolution, or the discrepancy must be included in $\zeta_k$.

\subsection{What contracts in the implemented and candidate observers}
\label{app:sampled-gradient}
The innovation bound controls error to truth along the realized gain
sequence. It does not omit the gain derivative from a pairwise Jacobian.
For scalar fixed data $z$ and inactive projection, put $e=z-h$ and
$\upsilon=e^2/\rho^2$. Differentiating the actual map gives
\[
\frac{d}{dh}\!\left(h+\frac{\gamma e}{1+\upsilon}\right)
=1-\frac{\gamma}{1+\upsilon}+\frac{2\gamma\upsilon}{(1+\upsilon)^2}
=1+\gamma\frac{\upsilon-1}{(1+\upsilon)^2}.
\]
The derivative of $(\upsilon-1)/(1+\upsilon)^2$ is $(3-\upsilon)/(1+\upsilon)^3$, proving the maximum
at $\upsilon=3$ in Eq.~\eqref{eq:innovation-counterexample}. When $z=f$, the
truth error is still multiplied by $1-\gamma/(1+\upsilon)$. This interior
counterexample rules out global Euclidean contraction of this map; it does
not rule out another metric on a restricted domain. Conversely, the legacy
map conditionally contracts under a common fixed observation stream but can
converge to the biased attenuated target. Neither property certifies
contraction of the robot--estimator feedback interconnection.

\begin{proposition}[A different sampled prediction-gradient candidate]
\label{prop:sampled-gradient}
Let $\Theta$ be closed and convex, $\Gamma\succ0$ fixed, and
$H_{I,k}=H_{I,k}^\top\succeq0$. Condition on a common exogenous data stream
$(H_{I,k},b_{I,k})$ and gains $h_k$ independent of the compared estimates.
The weighted prediction loss has gradient $H_{I,k}\theta-b_{I,k}$; its
projected negative-gradient step is
\begin{equation}
T_k(\theta)=\Pi_\Theta^{\Gamma^{-1}}
 \left[\theta+h_k\Gamma(b_{I,k}-H_{I,k}\theta)\right].
\label{eq:sampled-gradient}
\end{equation}
Define $S_{I,k}=\Gamma^{1/2}H_{I,k}\Gamma^{1/2}$ and
$q_{{\rm obs},k}=\max_i|1-h_k\lambda_i(S_{I,k})|$. Then
\begin{equation}
\|T_k(\theta_1)-T_k(\theta_2)\|_{\Gamma^{-1}}
\le q_{{\rm obs},k}\|\theta_1-\theta_2\|_{\Gamma^{-1}}.
\label{eq:sampled-gradient-contraction}
\end{equation}
For $S_{I,k}\succ0$, $0<h_k<2/\lambda_{\max}(S_{I,k})$ gives strict
conditional contraction. A full hold has $h_k=0$ and factor one.
\end{proposition}
\begin{proof}
Under $v=\Gamma^{-1/2}\theta$, metric projection is Euclidean projection
onto $\Gamma^{-1/2}\Theta$, hence nonexpansive even at active faces. The
unprojected difference is $(I-h_k S_{I,k})(v_1-v_2)$; its symmetric-matrix
operator norm is exactly $q_{{\rm obs},k}$. This proves the bound without
differentiating through projection. The step-size condition is equivalent
to $|1-h_k\lambda_i|<1$ for every positive eigenvalue.
\end{proof}
If $b_{I,k}=H_{I,k}\theta_k^*+\epsilon_{I,k}$ and
$\theta_k^*\in\Theta$, the same argument gives the truth-error bound
\begin{equation}
Z_{k+1}\le q_{{\rm obs},k}Z_k+
 h_k\|\Gamma^{1/2}\epsilon_{I,k}\|
 +\|\theta_{k+1}^*-\theta_k^*\|_{\Gamma^{-1}},
\label{eq:sampled-gradient-error}
\end{equation}
for $h_k\ge0$. Uniform exponential forgetting requires a uniform factor
below one or a proved windowed product condition. Rank-deficient information
precludes per-step strict contraction in every parameter direction. A FIR
command buffer alone is not an informative parameter dataset. Additive
$z_k=f_k+w_k$ with already known invertible sensitivity corresponds to
$H_I=I$, $b_I=z$, $\Gamma=I$; no new command excitation is then needed.
Unknown gains and offsets require separate regressor-rank conditions.

Equation~\eqref{eq:sampled-gradient} is an unevaluated candidate, not a
replacement interpretation of the reported attenuation or state-dependent
innovation runs. A gain selected from data quality or regressor energy is
admissible here only after conditioning on a common stream that does not
depend on the compared estimates. The actual robot stream changes when
commands change; the coupled analysis below addresses that separate issue.

\subsection{From an execution metric to a scheduled recovery tube}
\label{app:history-partition}

\paragraph{Execution state, command source, and observer state.}
The certified state $\xi_k$ contains the physical robot and necessary
low-level controller state, including actuator delays, filters, and reference
interpolation memory when required for a closed execution model. Manipulation
also requires the physical-state closure specified below. The VLA
supplies $a_k$; its observation history, network state, and policy
memory belong to the command source and are not part of $\xi_k$. The
observer's issued pre-fault command history, estimate, gains, and applied
snapshot belong to a separate state $\xi_k^{\rm obs}$. A received nominal
chunk may be treated as an external command schedule; if its interpolation
or phase is handled inside the execution controller, the required controller
memory must instead be represented in $\xi_k$. The chosen realization and
command schedule must be explicit.

\paragraph{Physical state closure during manipulation.}
\label{app:physical-closure}
If object pose, velocity, deformation, or evolving contact variables affect
execution, include their necessary state in $\xi_k$ and its nominal
comparator, together with the relevant contact mode. A discrete mode
indexes the vector field and metric; it is not a smooth Euclidean coordinate.
Fixed environment
geometry may be a common model parameter; prescribed environment motion may
be a declared common exogenous input. An object moved by the robot is an
endogenous physical state, not a common input merely because the commands
are shared. The metric and its domains must cover the selected physical
state and all allowed geometry, modes, and transitions. The actual and nominal
systems receive the same commands but can reach different object states and
contact sequences; agreement of those sequences is not assumed by notation.

Alternatively, a reduced servo state may omit these variables if their
remaining influence admits an explicit uniform transition-error bound
$\eta_k^{\rm contact}$ as specified below. This is a bound over admissible
corrected states, commands, omitted physical states, and contact events;
it need not be small or vanish under perfect actuator correction. Without
full state closure or such a bound, a servo/free-space certificate does not
extend to contact manipulation. Even a complete physical state does not
by itself establish contraction of its dynamics.

The sufficient closure condition for a $q_k$-only execution transition is:
at fixed $\xi_k$ and supplied command, any two observer/adapter states
producing the same realized effective mismatch $q_k$ must produce the same
next execution state. With no separate transition disturbance, $q_k=0$
then recovers the nominal execution map. A corrected-command queue may be
outside $\xi_k$ only if its entire effect is mediated by the command
actually executed; plant or controller delays cannot simply be discarded.
Under this closure, condition on the realized supplied commands and write
\begin{equation}
\xi_{k+1}^{\rm nom}=F_{\mathrm{exec},k}^0(\xi_{k}^{\rm nom},a_k),\qquad
\xi_{k+1}=F_{\mathrm{exec},k}(\xi_k,a_k,q_k).
\label{eq:app-closed-loop}
\end{equation}
The reference is the nominal physical response to these same commands.
It does not rerun the VLA on hypothetical healthy observations. The commands
may be generated using the actual robot's observations: the result is then
a pathwise execution comparison, valid while those commands and the states
remain in the assumed admissible sets. It does not establish stability of
the feedback law that generated the commands. Computing $\xi_{k}^{\rm nom}$ online
is unnecessary for the evaluated prediction-only laws.

For binary $D$ and box projection,
\begin{equation}
\|q_k\|\le U_k^\perp+E_k^{D,\rm app}+\bar\zeta_k,
\qquad U_k^\perp=\|(I-D)f_k\|.
\label{eq:app-mask}
\end{equation}
The supplied nominal interface action $a_k$ is the one used to define the
actual mismatch; it is not a newly queried healthy-policy action. For
general $D$ or a coupled projection use instead
$\|q_k\|\le U_k^\perp+\|D\|B_k^{\rm app}+\bar\zeta_k$.
Native-decoder error and command clipping belong in $\zeta_k$; any other
separate transition disturbance belongs in $\eta_k$ below. Each physical
effect is counted once. In particular, bounded parameter projection does
not imply command feasibility or repair of unmatched joint faults.

\paragraph{Additional history and contact effects.}
If a plant or low-level controller needs a history not determined by
$(\xi_k,a_k,q_k)$, retain the required memory in $\xi_k$ and its
nominal comparator, or use an effective-input history with an exact
realization. A VLA policy that reads raw issued-action history changes the
supplied command stream; this dependence alone does not place its memory
inside the primary execution certificate. If a raw-command or corrected-queue
update has an additional direct effect on the certified execution state,
write its explicit map as $\mathcal F_{\mathrm{exec},k}(\xi,a,q,h)$ with
$\mathcal F_{\mathrm{exec},k}(\xi,a,q,0)=F_{\mathrm{exec},k}(\xi,a,q)$.
Require, uniformly on the certified domain and over allowed commands,
\begin{equation}
\begin{aligned}
d_{k+1}(\mathcal F_{\mathrm{exec},k}(\xi,a,q,h),F_{\mathrm{exec},k}(\xi,a,q))
&\le L_k^h\|h\|,\\
\|h_k\|&\le\bar h_k,\\
d_{k+1}(\xi_{k+1},\mathcal F_{\mathrm{exec},k}(\xi_k,a_k,q_k,h_k))
&\le\eta_k^{\rm contact},\\
\eta_k=\eta_k^{\rm base}+\eta_k^{\rm hist}+\eta_k^{\rm contact},\qquad
\eta_k^{\rm hist}&=L_k^h\bar h_k.
\end{aligned}
\label{eq:app-history-forcing}
\end{equation}
The third line bounds the remaining omitted contact/environment effect on
the reduced-state transition; it is zero when that effect is already
represented exactly. Its validity must include the relevant event timing
and resets. Here $\eta_k^{\rm base}$ bounds other base-map disturbance,
excluding the displayed history/contact effects and realization error
already included in $q_k$. The triangle inequality gives
\[
\begin{aligned}
d_{k+1}(\xi_{k+1},F_{\mathrm{exec},k}^0(\xi_k,a_k))
&\le L_k\|q_k\|+\eta_k^{\rm base}+L_k^h\|h_k\|+
\eta_k^{\rm contact}\\
&\le L_k\|q_k\|+\eta_k.
\end{aligned}
\]
These are required bounds, not definitions that make omitted effects small
or make them vanish at $q_k=0$. In a reduced model, the residual bound is
required for every admissible hidden-state realization. A different
execution-state realization requires its own distance-contraction and
transition check. The tube proof
uses total forcing $\eta_k$, with each physical effect counted once.

For example, take $p_{k+1}=.8p_k+.2(u_k+f)$, supplied nominal $a_k=0$, a
constant nonzero fault $f$, and exact correction $u_k=-\widehat f_k=-f$.
Then $q_k=0$ and physical motion matches $p_{k+1}^{\rm nom}=.8p_k^{\rm nom}$ from equal
initial positions. Nevertheless, a raw-command memory $v_{k+1}=u_k$ equals
$-f$, whereas $v_{k+1}^{\rm nom}=0$. Starting with equal $(p_0,v_0)$, their next
augmented distance is $|f|$, so a zero-forcing $q$-only tube for $(p,v)$
fails even though its nominal Euclidean map contracts by $.8$. For the
physical execution state $\xi=p$, the $q$-only model is exact; $v$ can stay
in the observer. If $(p,v)$ is certified instead, its exact augmented map
$\mathcal F((p,v),q,h)=(.8p+.2q,h)$ with $h=u-a$ has $L=.2$, $L^h=1$,
and $\bar h=|f|$ under exact correction. The conditional tube proof is
unchanged; its state and forcing must represent the chosen subsystem.

\paragraph{Direct contraction and smooth sufficient conditions.}
Theorem~\ref{thm:tube} assumes the distance contraction in
Eq.~\eqref{eq:sampled-contraction} directly. For a smooth-mode sufficient
construction, let $M_{c,k}(\xi)$ be smooth, independent of the supplied
command and adapted estimate, and satisfy
$0<\underline m I\preceq M_{c,k}\preceq\overline m I$ on the stated domains;
let $d_k$ be its intrinsic Riemannian distance. With Euclidean
$\delta_{\rm phys}$ these bounds imply Eq.~\eqref{eq:distance-equivalence}
with $c_-=\sqrt{\underline m}$ and $c_+=\sqrt{\overline m}$. Suppose
$F_{\mathrm{exec},k}^0(\cdot,a)$ is $C^1$ and, uniformly over allowed $a$,
\begin{equation}
J_k^\top M_{c,k+1}(F_{\mathrm{exec},k}^0(\xi,a))J_k
\preceq(\lambda_k^{\rm exec})^2M_{c,k}(\xi),\qquad
J_k=\partial_\xi F_{\mathrm{exec},k}^0(\xi,a),
\label{eq:sampled-metric}
\end{equation}
with $a$ held fixed. If admissible connecting curves have admissible nominal
images, this inequality contracts each image-curve length by
$\lambda_k^{\rm exec}$; taking the infimum proves
Eq.~\eqref{eq:sampled-contraction}. Thus Eq.~\eqref{eq:sampled-metric} is
sufficient in a smooth mode but is not required when
Eq.~\eqref{eq:sampled-contraction} is established directly. A geodesically
convex neighborhood with its relevant images contained in the next domain is
sufficient for this construction; no VLA Jacobian appears.

Separately, Theorem~\ref{thm:tube} assumes the realized-transition bound in
Eq.~\eqref{eq:sampled-transition}. Under this Riemannian construction, if
$F_{\mathrm{exec},k}$ is differentiable in its actual physical input, a
sufficient input-sensitivity condition is
\[
\left\|M_{c,k+1}^{1/2}(F_{\mathrm{exec},k}(\xi,a,v))
D_qF_{\mathrm{exec},k}(\xi,a,v)\right\|\le L_k
\quad\text{for every }v=sq,\ 0\le s\le1.
\]
Integrating the length of $s\mapsto F_{\mathrm{exec},k}(\xi,a,sq)$, then
adding $d_{k+1}(F_{\mathrm{exec},k}(\xi,a,0),F_{\mathrm{exec},k}^0(\xi,a))
\le\eta_k$, proves
\begin{equation}
d_{k+1}(F_{\mathrm{exec},k}(\xi,a,q),F_{\mathrm{exec},k}^0(\xi,a))
\le L_k\|q\|+\eta_k.
\label{eq:app-input-assumption}
\end{equation}
For the explicit decomposition in \eqref{eq:app-history-forcing}, the
base-map argument uses $\eta_k^{\rm base}$ instead; adding its history
and contact bounds once gives the same inequality for the realized next
state and hence Eq.~\eqref{eq:sampled-transition}. The entire input path must
lie in the domain. For a nonsmooth map, establish
Eq.~\eqref{eq:sampled-transition} directly. The condition concerns the
realized physical input, not merely the pre-decoder correction parameter.

\paragraph{Proof and quantitative form of Theorem~\ref{thm:tube}.}
Insert $F_{\mathrm{exec},k}^0(\xi_k,a_k)$ between the two next states and apply the triangle
inequality, Eqs.~\eqref{eq:sampled-contraction} and
\eqref{eq:sampled-transition}, and the masked mismatch bound. Thus
$X_{k+1}\le \lambda_k^{\rm exec}X_k+g_k$, with $X_k=d_k(\xi_k,\xi_{k}^{\rm nom})$, where
\begin{equation}
g_k=L_k(U_k^\perp+B_k^{D,\rm app}+\bar\zeta_k)+\eta_k,
\qquad R_{k+1}=\lambda_k^{\rm exec}R_k+g_k.
\label{eq:app-tube-recursion}
\end{equation}
Induction proves inclusion from $R_0\ge d_0(\xi_0,\xi_{0}^{\rm nom})$ as long as the
radius-$R_k$ balls remain in the domains. Expansion gives
\begin{equation}
R_k=\left(\prod_{i=0}^{k-1}\lambda_i^{\rm exec}\right)R_0+
\sum_{j=0}^{k-1}\left(\prod_{i=j+1}^{k-1}\lambda_i^{\rm exec}\right)g_j.
\label{eq:app-tube-product}
\end{equation}
Arbitrary finite nonnegative amplification factors give only this
finite-horizon statement. If every subinterval satisfies
$\prod_{i=j}^{k-1}\lambda_i^{\rm exec}\le C\varrho^{k-j}$, $C\ge1$, $0<\varrho<1$, then
\begin{equation}
R_k\le C\varrho^kR_0+C\sum_{j=0}^{k-1}\varrho^{k-1-j}g_j
\le C\varrho^kR_0+\frac{C\bar g}{1-\varrho}(1-\varrho^k)
\label{eq:app-product-recovery}
\end{equation}
when $g_j\le\bar g$. This allows transient expansion; a uniform
$\lambda_k^{\rm exec}\le a<1$ is the special case $C=1,\varrho=a$.
If $0\le a,q<1$, $\lambda_k^{\rm exec}\le a$, $L_k\le L$, $U_k^\perp\le U^\perp$,
$\bar\zeta_k\le\bar\zeta$, $\eta_k\le\eta$, and
$B_k^{D,\rm app}\le q^k E_0^D+\bar E$, then
\begin{equation}
R_k\le a^kR_0+LE_0^D\mathcal C_k(a,q)+
\frac{L(U^\perp+\bar E+\bar\zeta)+\eta}{1-a}(1-a^k),
\label{eq:app-tube-geometric}
\end{equation}
where $\mathcal C_k(a,q)=\sum_{j=0}^{k-1}a^{k-1-j}q^j$ equals
$(a^k-q^k)/(a-q)$ if $a\ne q$, and $ka^{k-1}$ if $a=q$ for $k\ge1$;
$\mathcal C_0=0$. Bounds on the applied error must already include its actual
schedule; estimate convergence alone does not imply this geometric form.

\paragraph{Hybrid execution, chunks, and reference jumps.}
At a step containing an impact, reset, or mode change,
Eq.~\eqref{eq:sampled-contraction} must hold for the complete nominal sampled
transition; a within-mode Jacobian check alone is insufficient.
If the domains form a hybrid union, $d_k$ must be defined for the cross-mode
state pairs used by the comparison, or the certified event structure must
exclude those pairs.
The physical certificate includes relevant controller memory, actuator
delays, and contact modes. Within a smooth mode, flow inequalities apply
only on its certified domain. Impulsive contact forces are not ordinary
bounded continuous disturbances: impacts require state/metric jump bounds,
or a verified sampled transition bound that includes the impact.
State-dependent event timing requires a valid saltation or incremental
event-map bound; a reset Jacobian alone can omit event-time sensitivity.
The compared executions need not impact at the same time. A common-mode
flow estimate alone does not compare differing mode sequences; either
justify the required hybrid comparison, restrict the admissible events,
or bound the resulting sampled transition discrepancy explicitly.
Uniform bounds can fail near grazing events or changes in contact sequence. A chunk-boundary metric may be appropriate
when buffer shifts obstruct one-step contraction, but intermediate physical
motion requires its own bounds. Boundary stability does not certify
intermediate constraints.

Both physical systems receive the same command changes. In physical-state
coordinates their responses need not jump when the command changes; no
command derivative enters a command-independent metric. A commanded reset
of low-level controller state must nevertheless be included in the
transition. If tracking coordinates are instead measured directly from a
commanded reference that jumps, retain the resulting coordinate/reset jump
and its metric gain or additive forcing, or use a feasible interpolated
reference. Common commands cannot justify omitting a reset that the chosen
coordinates actually undergo. This conditional pathwise result is neither
a distributional success theorem nor a certificate of the VLA command
source's feedback stability.

\paragraph{Optional conditional corollary: task preservation.}
\label{app:task-extension}
Let $\mathcal S$ be the successful physical trajectories over horizon $N$,
including any specified command-dependent constraints. Suppose the nominal
execution trajectory $\xi_{0:N}^{\rm nom}$ for the supplied command sequence is
successful and has positive radii $\mu_k^{\rm task}$ such that every
admissible trajectory with $d_k(\xi_k,\xi_{k}^{\rm nom})<\mu_k^{\rm task}$ at
all steps belongs to $\mathcal S$. If the execution-tube assumptions hold
and $R_k<\mu_k^{\rm task}$ throughout, the realized execution succeeds.
This follows immediately from tube inclusion and the assumed robustness
neighborhood. The neighborhood must cover all path and terminal conditions
and intermediate chunk motion; constraints on omitted variables require
separate bounds. Neither success of the supplied nominal response nor
its whole-task margin follows from contraction.

When comparing two adapters, their observation feedback may produce
different VLA command streams. Fixed-command replay gives the common-input
comparison directly; otherwise each adapter first has its own nominal
physical response. In a common certified metric, the triangle inequality
then bounds their physical separation by the sum of their two execution
radii and the distance between those nominal responses. That last
command-divergence term requires separate evidence and cannot be set to
zero because the VLA weights are frozen.

An independently rerun healthy VLA generally supplies different commands.
Comparing against that rollout requires a separate bound on the divergence
between its nominal physical response and $\xi_{k}^{\rm nom}$. For example, if such
a bound is $d_k(\xi_{k}^{\rm nom},\xi_k^h)\le C_k^{\rm cmd}$, the triangle
inequality gives $d_k(\xi_k,\xi_k^h)\le R_k+C_k^{\rm cmd}$; a task
margin around $\xi_k^h$ must exceed this sum. A bound on command differences
together with a valid execution input-sensitivity bound can provide
$C_k^{\rm cmd}$, but it must be established separately. Neither a VLA
metric nor this outer-loop comparison is needed for the primary execution
tube. The reported experiments establish no task margins. A uniform physical
execution metric was measured directly by same-command replay rather than
inferred from fitted gains or sampled Jacobians, and exists on the GR1 arm
but not on the Panda under operational-space control
(Appendix~\ref{app:measured-certificates}).

\subsection{Closing the loop when predictor error depends on trajectory error}
\label{app:smallgain}

An independently bounded $\epsilon_k$ can be a strong assumption: a correction
changes the visited physical states and commands, and predictor error can
grow with execution deviation and estimation error. The assumed error bound
must hold uniformly over admissible supplied commands; it does not follow
from a healthy-data fit or from the physical metric alone. The following result gives a sufficient
condition for closing this feedback loop. It applies when every physical step updates and applies the current estimate,
with a uniformly positive active step size. Observation attenuation is included
by adding its forcing to the constant observation-error term below. Arbitrary
holds and delayed application instead use the scheduled tube.

\begin{proposition}[A conditional small-gain bound]
\label{prop:smallgain}
Let $X_k=d_k(\xi_k,\xi_{k}^{\rm nom})$ be the same-command execution distance and write
$E_k=E_k^D$ within this subsection. Suppose on
the relevant local domain
\begin{align}
 X_{k+1}&\le\lambda X_k+b E_k+d_0,\nonumber\\
 E_{k+1}&\le[1-\alpha_k(1-k_E)]E_k+
       \alpha_k(\epsilon_0+k_X X_k)+\nu,
 \label{eq:app-smallgain-recursions}
\end{align}
where $0\le\lambda<1$, $0\le k_E<1$, $b,k_X,\epsilon_0,\nu,d_0\ge0$ and
$0<\underline\alpha\le\alpha_k\le\overline\alpha\le1$.
The first inequality follows from the tube assumptions with
$b=L$ and $d_0=L(U^\perp+\overline\zeta)+\eta$;
equivalently, an adapter discrepancy can be absorbed into
$\eta_k$ as $L_k\overline\zeta_k$, but must not be counted twice.
For innovation adaptation the second follows from
$\|Dw_k\|\le\epsilon_0+k_E E_k+k_X X_k$. For the observation-attenuated
law take $\alpha_k=\gamma$ and replace $\epsilon_0$ throughout this
proposition by $\epsilon_0+\bar a_{\rm att}$. If
\begin{equation}
 k_Xb<(1-k_E)(1-\lambda),
 \label{eq:app-smallgain-condition}
\end{equation}
choose $s>0$ with
$b/(1-\lambda)<s<(1-k_E)/k_X$, with the upper bound
interpreted as $+\infty$ when $k_X=0$.
Then $W_k=\max\{X_k,s E_k\}$ satisfies
\begin{equation}
 W_k\le\varrho^k W_0+\frac{d_*}{1-\varrho}(1-\varrho^k),
 \label{eq:app-smallgain-bound}
\end{equation}
where
\begin{equation}
 \varrho=\max\{\lambda+b/s,\ 1-\underline\alpha(1-k_E-sk_X)\}<1,
 \qquad d_*=\max\{d_0,\ s(\overline\alpha\epsilon_0+\nu)\}.
 \label{eq:app-smallgain-constants}
\end{equation}
The bound holds inductively as long as its implied trajectory and estimator
neighborhoods remain inside the domain of the assumptions.
\end{proposition}
\begin{proof}
The strict small-gain inequality makes the interval for $s$ nonempty. Since
$X_k\le W_k$ and $E_k\le W_k/s$, the first recursion gives
$X_{k+1}\le(\lambda+b/s)W_k+d_0$. The second gives
\begin{align*}
 s E_{k+1}
 &\le[1-\alpha_k(1-k_E)+sk_X\alpha_k]W_k+s(\alpha_k\epsilon_0+\nu)\\
 &\le[1-\underline\alpha(1-k_E-sk_X)]W_k
       +s(\overline\alpha\epsilon_0+\nu).
\end{align*}
Both coefficients are strictly below one. Taking the maximum and expanding
$W_{k+1}\le\varrho W_k+d_*$ proves the claim.
\end{proof}

For a constant step size, the homogeneous comparison matrix is
\[
 A_\alpha=\begin{pmatrix}\lambda&b\\
             \alpha k_X&1-\alpha(1-k_E)\end{pmatrix}.
\]
The weighted-maximum proof also covers the stated time-varying family.
The result is a stability bound for the coupled error comparison, not a claim
of global pairwise contraction of the nonlinear estimator. It explains why a
well-conditioned sensitivity and small predictor error growth can matter even
when healthy prediction fits appear good. The constants $k_E,k_X,\lambda$ were measured on the deployed plant
(Appendix~\ref{app:measured-certificates}): the coupling terms are small
($k_E=.012$, $k_X=.009$, $\epsilon_0=.0008$), and the condition nonetheless
fails because $a=1-\lambda=9\cdot10^{-5}$. The obstruction is the absent
contraction, not the feedback coupling this subsection bounds.

\subsection{Invariant recovery regions and admissible startup error}
\label{app:discrete-region}
\paragraph{Proof of Corollary~\ref{thm:discrete-region}.}
Throughout this subsection, $E_k$ abbreviates $E_k^D$, and $\nu$ bounds
masked drift.
In the active-every-step case with current application, the two estimator
recursions imply the common comparison
\[
E_{k+1}\le[1-\beta_k c]E_k+\beta_k(k_XX_k+v),
\]
where $\beta_k=\gamma$ for observation attenuation and $\beta_k=\alpha_k$
for innovation adaptation. For the first case use
$(1-s_k)\|Df_k\|\le\bar a_{\rm att}$ and $s_k\le1$; for the second
use $\nu\le\alpha_k\nu/\underline\alpha$.
The coefficients multiplying $X_k,E_k$ are nonnegative because
$0\le k_E<1$ and $0<\beta_k\le1$. At the upper corner of the
rectangle, its second inequality gives
$\beta_k(k_X\bar X+v)\le\beta_k c\bar E$. Hence
$E_{k+1}\le\bar E$.
The first inequality gives $X_{k+1}\le\bar X$. Induction from the stated
initial conditions proves invariance, and the assumed domain containment
makes each application valid. Solving the two corner equalities gives
the equality corner
\begin{equation}
X_* =\frac{cd_0+bv}{\Delta},\qquad
E_* =\frac{k_Xd_0+av}{\Delta}.
\label{eq:discrete-corner}
\end{equation}
A task-preservation conclusion additionally
requires the successful nominal response and whole-trajectory neighborhood
in Appendix~\ref{app:task-extension}. \hfill$\square$

For a chosen $\bar X<\mu_{\rm exec}$, where $\mu_{\rm exec}$ is a
valid execution-domain tube radius, and $b>0$, a sufficient startup budget
given valid parameter-error and domain bounds is
\begin{equation}
bv+c d_0\le\Delta\bar X,
\quad bE_0+d_0\le a\bar X,\quad X_0\le\bar X,
\label{eq:discrete-budget}
\end{equation}
with estimator radius chosen from the nonempty interval
\begin{equation}
\max\{E_0,(k_X\bar X+v)/c\}\le\bar E\le(a\bar X-d_0)/b.
\label{eq:discrete-radius-choice}
\end{equation}
The second condition in Eq.~\eqref{eq:discrete-budget} is essential:
the steady-state budget alone does not accommodate an arbitrarily large
fault-onset estimation error. For $b=0$, use the two original rectangle
inequalities directly. A common scaling of $(X_*,E_*)$ by a factor at least
one preserves both inequalities. It can cover arbitrary finite initial errors
when both corner components are positive. If either is zero, use the original
rectangle inequalities or the product bound directly.

For completeness, the primary observation-attenuated law has the explicit
comparison below. Suppose
$\bar a_{\rm att}\ge(1-s_k)\|Df_k\|$ and
$\|Dw_k\|\le\epsilon_0+k_EE_k+k_XX_k$. Since $0\le s_k\le1$,
the masked counterpart of Eq.~\eqref{eq:app-legacy-bound} implies
\begin{equation}
E_{k+1}\le[1-\gamma(1-k_E)]E_k+
\gamma(\epsilon_0+\bar a_{\rm att}+k_XX_k)+\nu.
\label{eq:discrete-legacy-region}
\end{equation}
The same rectangle proof applies with
$v=\epsilon_0+\bar a_{\rm att}+\nu/\gamma$.
It is conservative; retaining $s_k$ in the original recursion can be
tighter. The attenuation floor cannot be removed when this is the executed law.

\paragraph{An exact step-history contribution.}
For an otherwise constant fault jumping by $\Delta f$ at $k_0$, the
history part of Eq.~\eqref{eq:app-observation-error} at $j=k-k_0\ge0$ is
\begin{equation}
w_{k_0+j}^{\rm hist}=-S^{-1}\sum_{\ell=j+1}^{K}G_\ell\Delta f.
\label{eq:discrete-onset}
\end{equation}
For each term with $\ell>j$, the delayed fault precedes the jump and
$f_{k-\ell}-f_k=-\Delta f$; every other term is zero. This proves the
identity and its disappearance after $K$ steps. Substituting its time-dependent
norm into the estimator and trajectory recursions accounts for onset without
treating the jump as a persistent drift rate.

\paragraph{Rate, noise, and drift are separate quantities.}
For fixed $s$, constant noiseless observations, and inactive projection,
the legacy transient is
$\widehat f_k-sf=(1-\gamma)^k(\widehat f_0-sf)$.
With $\gamma=.08$, 36 updates reduce that transient by at least 95\%;
this concerns the attenuated target. For a constant-gain innovation observer
with independent zero-mean scalar measurement noise of variance $\sigma_w^2$,
inactive projection, and constant fault, the stationary error variance is
$\alpha\sigma_w^2/(2-\alpha)$, obtained by solving
$v=(1-\alpha)^2v+\alpha^2\sigma_w^2$. This calculation is illustrative:
correlated FIR residuals, normalization, and gates require their own
stochastic analysis. The deterministic drift bound remains $\nu/\alpha$.

\subsection{Why identifying first and then holding changes the bound}

Suppose an estimate obtained before an evaluation episode is held at
$\widehat f^H$, with $\|D(f_0-\widehat f^H)\|\le E_H$. If the persistent
fault can drift in selected coordinates by at most $\nu$ per control step, then
\begin{equation}
 \|D(f_k-\widehat f^H)\|\le E_H+k\nu,
 \label{eq:app-held}
\end{equation}
by telescoping. Substituting this quantity for $B_k^{D,\rm app}$ in
\eqref{eq:app-tube-recursion} yields the held-estimate guarantee. For a constant
fault and a valid initialization, holding removes the online estimation
transient from the task episode, but retains calibration error; for a changing
fault it accumulates drift. A prior identification episode also changes the
information available at evaluation time. Its benefit cannot be attributed
solely to removing within-episode variance or compared as if both protocols
started with zero fault information.

\subsection{Gain faults and adapter dimension}
\label{app:gain-authority}

For a diagonal per-channel FIR and a gain $g_i$ constant over the FIR window,
the correct scalar regression is
\begin{equation}
 e_{i,k}^{\rm res}=\beta_i\phi_{i,k}+\nu_{i,k},\quad
 \beta_i=g_i-1,\quad
 \phi_{i,k}=\sum_{\ell=0}^{K}h_{i,\ell}u_{i,k-\ell}
           =\widehat y_{i,k}-c_i.
 \label{eq:app-gain-regressor}
\end{equation}
The equality is exact only for the stated diagonal, correctly calibrated
plant; coupling and time-varying gains introduce model error. In particular,
gain drift contributes
$\sum_\ell h_{i,\ell}(g_{i,k-\ell}-g_{i,k})u_{i,k-\ell}$.
Noiseless constant-gain identification requires
$\sum_k\phi_{i,k}^2>0$; robustness depends on its magnitude. With an unknown
constant intercept, the two-column regression $[\phi_{i,k},1]$ has full rank
only when $\phi_{i,k}$ varies. Command quantile ranges alone establish neither
of these conditions.

The inverse-gain adapter $u_i=a_i/\widehat g_i$, with
$\widehat g_i\ge g_{\min}>0$, gives
\begin{equation}
 q_i=\left(\frac{g_i}{\widehat g_i}-1\right)a_i
     =\frac{g_i-\widehat g_i}{\widehat g_i}a_i,
 \qquad
 |q_i|\le\frac{|a_i|}{g_{\min}}|g_i-\widehat g_i|.
 \label{eq:app-gain-mismatch}
\end{equation}
This is the physical mismatch to insert in the trajectory tube. The derivative
of the \emph{command} with respect to its estimate is
$-a_i/\widehat g_i^2$; the derivative of the \emph{executed command} is
$-g_i a_i/\widehat g_i^2$. They have different amplification factors.
Input limits must still permit the required inverse command.

\begin{proposition}[Local correction authority and dimension]
\label{prop:directions}
Suppose a family of constant faults is $f=V\theta$ and an additive
adapter can generate $B\vartheta$ with fixed $B\in\mathbb R^{d\times p}$. Exact
action-level compensation of the whole family by unrestricted adapter
coordinates is possible if and only if
\begin{equation}
 \operatorname{Range}(V)\subseteq\operatorname{Range}(B),
 \qquad\text{which implies}\qquad
 p\ge\operatorname{rank}(V).
 \label{eq:app-rank}
\end{equation}
\end{proposition}
\begin{proof}
For necessity, apply the matching requirement to each coordinate basis vector
$\theta$; every column of $V$ must then lie in the range of $B$.
For sufficiency, choose $K$ with $BK=V$ and set $\vartheta=K\theta$
(with the compensation sign in the commanded action).
\end{proof}
With bounded adapter coordinates, the required solutions must also be feasible
inside those bounds; range inclusion alone does not ensure that. If only a local
physical effect matters, the corresponding criterion uses
$\operatorname{Range}(A V)\subseteq\operatorname{Range}(A B)$ for the
linearized action-to-effect map $A$. For one shared parameter vector across
multiple times, the time-indexed maps must be stacked before applying this
range criterion. A six-channel interface spans arbitrary six-dimensional
offsets, whereas a known uniform offset $f=\alpha\mathbf 1$ is a one-dimensional
family. The experiments do not establish that six is a universal minimum.
A joint lock can violate local matching or command feasibility; it does not
prove that another policy or a different path through the remaining joints
could never complete the task.
